\documentclass[12pt]{article}

\usepackage{graphicx}
\usepackage{longtable}

\newcommand{\DocumentName}{Software Test Documentation}
\newcommand{\Title}{\DocumentName \\ \hfill}
\newcommand{\Author}{Youssef Boudiba}
\newcommand{\Version}{1.0}

\begin{document}

\begin{titlepage}

\begin{center}
\includegraphics[width=250pt,angle=90]{../../templates/img/logopdf.pdf}
\vfill
{\Huge \Title}
\hfill
\vfill
\vfill
\end{center}

\begin{minipage}[t]{1\textwidth}
\begin{flushleft}
\includegraphics[width=100pt]{../../templates/img/VUB_logo_compact.jpg}
\end{flushleft}
\end{minipage}

\begin{minipage}[t]{1\textwidth}
\begin{flushright}
\Author \\
\Version \\
\end{flushright}
\end{minipage}

\end{titlepage}
\section{Introduction}\label{introduction}

The software test documentation is based on the IEEE standard for
software test documentation (829-1998). Since we are working on a small
project, we will not cover every topic proposed by the standard, but
will instead focus on those who are more relevant to our project.

This document consists of 3 parts:

\textbf{\emph{1. Test specification:}}

Consists of describing (major parts of) the software test plan for our
project. By defining a test plan we will be able to cover the management
of the tests, how we are going to run those tests, what we are going to
tests, etc.

\textbf{\emph{2. Test monitoring:}}

We will keep tracks of the tests through a test log and report incidents
via an incident report in the second part of the document.

\textbf{\emph{3. Test report:}}

Last but not least is the test report. This will be an overview
(summary) of some important tests as well as a conclusion of the
software test documentation.

\textbf{Record of revision}

\begin{longtable}[c]{@{}lll@{}}
\hline\noalign{\medskip}
Version & Date & Description
\\\noalign{\medskip}
\hline\noalign{\medskip}
\textbf{0.1} & 8/11/2013 & Initial draft
\\\noalign{\medskip}
\textbf{0.2} & 10/11/2013 & Added additional content
\\\noalign{\medskip}
\textbf{0.3} & 8/12/2013 & Removing unnecessary content and updating topics
\\\noalign{\medskip}
\textbf{1.0} & 12/12/2013 & Ready for delivery: iteration 1
\\\noalign{\medskip}
\hline
\end{longtable}

\textbf{References}

\begin{enumerate}
\def\labelenumi{\arabic{enumi}.}
\itemsep1pt\parskip0pt\parsep0pt
\item
  \textbf{\emph{Software project management plan (SPMP):}}
\end{enumerate}

https://github.com/se1-1314/xiast-docs/blob/master/management/project/spmp.pdf

\begin{enumerate}
\def\labelenumi{\arabic{enumi}.}
\setcounter{enumi}{1}
\itemsep1pt\parskip0pt\parsep0pt
\item
  \textbf{\emph{Software requirements specifications (SRS):}}
\end{enumerate}

https://github.com/se1-1314/xiast-docs/blob/master/management/requirements/srs.pdf

\section{Test Specification}\label{test-specification}

This section describes the software test plan.

\subsection{Approach (Testing
strategy)}\label{approach-testing-strategy}

Since we are following an iterative and incremental development model
for our project (check SPMP for more info), our test strategy will
consist of a series of different tests. These tests are specified below
and are done by the whole group since the absence of a testing team.

\begin{enumerate}
\def\labelenumi{\arabic{enumi}.}
\item
  \textbf{Unit testing}

  Every small amount of source code (or low level software module) that
  cannot be divided in smaller part is called a unit. Whenever a
  developer has finished working on such a unit, it will be tested to
  check whether or not it is conform to its design and requirements.
\item
  \textbf{Integration testing}

  Integration testing combine and test low level software modules that
  have been unit tested. This test assures us that the combined higher
  level module still works and that it is conform to its design and
  requirements.
\item
  \textbf{System testing}

  System testing tests and combines all integrated (higher level)
  software component. The test will be performed on the whole system in
  order to find deviations from the specified requirments.
\end{enumerate}

\subsection{Features (not) to be
tested}\label{features-not-to-be-tested}

Since we'll most likely follow the test driven development cycle (TDD),
every unit will have its test case. If changing a unit's content doesn't
cause any test to fail, then the unit will be consider useless and must
be deleted (unless a test is missing to support its existence).

\subsection{Pass/fail criteria}\label{passfail-criteria}

A test pass if no error occurs and if it's conform to its design and
requirements. Whenever an error occurs whilst running a test, the
responsible developer will open an issue on gitHub and try to find a
solution (with whole team if necessary). When a solution is found, a
test incident report must be generated to keep track of the occurred
error (see 3. test execution for more info).

\subsection{Testing tools}\label{testing-tools}

\paragraph{clojure.test}\label{clojure.test}

The clojure library itself provide testing tools for testing the code
and can be found at 
http://richhickey.github.io/clojure/clojure.test-api.html. Since
no testing framework has been decided yet we'll use this library in the
meantime.

\section{Test Monitoring}\label{test-monitoring}

Both the treatment of errors as keeping track of tests results is
described below in this section. We will do this by generating two kinds
of documents:

\subsection{Test log}\label{test-log}

A test log keeps track of what happened during a test and must contain:

\begin{enumerate}
\def\labelenumi{\arabic{enumi}.}
\item
  \textbf{\emph{Execution description}}
\item
  \textbf{\emph{Results}}
\end{enumerate}

We will be using gitHub to store the test log under a specific
directory.

\subsection{Test incident report}\label{test-incident-report}

We will keep track of incidents (bugs, unexpected crashes, unexpected
results of a test, etc.) by using the issue tool on gitHub and
generating a report whenever and issue has been resolved. Such a report
consists of:

\begin{enumerate}
\def\labelenumi{\arabic{enumi}.}
\item
  \textbf{\emph{Problem identifying}}

  \begin{enumerate}
  \def\labelenumii{\arabic{enumii}.}
  \item
    \textbf{\emph{Incident description}}
  \item
    \textbf{\emph{Impact on system components}}
  \end{enumerate}
\item
  \textbf{\emph{Problem solving}}

  \begin{enumerate}
  \def\labelenumii{\arabic{enumii}.}
  \setcounter{enumii}{2}
  \item
    \textbf{\emph{Priority}}
  \item
    \textbf{\emph{Actions taken}}
  \item
    \textbf{\emph{Solution}}
  \end{enumerate}
\end{enumerate}

We will be using gitHub to store the test incident reports under a
specific directory.

\section{Test Report}\label{test-report}

This section contains the result and evaluation of important tests as
well as a conclusion to the software test document.

\end{document}
